\documentclass[prl,aps,twocolumn,notitlepage,superscriptaddress,showpacs,nofootinbib]{revtex4-2}

%\documentclass[prl,aps,twocolumn,notitlepage,superscriptaddress,showpacs,nofootinbib]{revtex4-2}
%\documentclass[showpacs,amsmath,amssymb,twocolumn,prl,superscriptaddress]{revtex4-1}
%\usepackage{graphicx} % Required for inserting images

\input{package}
\newcommand{\red}[1]{{\color{red} #1}}

\hypersetup{
	colorlinks = true,
	urlcolor = {blue},
	citecolor = {blue},
	linkcolor= {blue}
}

\RequirePackage[switch,columnwise]{lineno}
\makeatletter
\newif\ifLNrightcolumn
\let\LNoriggotNumberedPage\gotNumberedPage
\def\gotNumberedPage#1#2#3#4{%
  \ifnum#3=2\relax
    \global\LNrightcolumntrue
  \else
    \global\LNrightcolumnfalse
  \fi
  \LNoriggotNumberedPage{#1}{#2}{#3}{#4}}
\def\makeLineNumberRight{%
  \linenumberfont\hskip\linenumbersep
  \hskip\hsize
  \hb@xt@\linenumberwidth{\hss\LineNumber}\hss}
\let\makeLineNumberOdd\makeLineNumberLeft
\let\makeLineNumberEven\makeLineNumberRight
\def\makePagewiseLineNumber{%
  \logtheLineNumber\getLineNumber
  \ifLNrightcolumn
    \makeLineNumberRight
  \else
    \makeLineNumberLeft
  \fi}
\makeatother

\begin{document}
\title{Symmetry speeds up quantum measurements}

\author{Hu Chen}
\altaffiliation{These authors contributed equally to this work.}
\affiliation{Southern University of Science and Technology, Shenzhen, 518055, China}
\affiliation{International Quantum Academy, Shenzhen, 518048, China}

\author{Bujiao Wu}
\altaffiliation{These authors contributed equally to this work.}
\affiliation{International Quantum Academy, Shenzhen, 518048, China}

\author{Qian-Xi Zhang}
\affiliation{Southern University of Science and Technology, Shenzhen, 518055, China}
\affiliation{International Quantum Academy, Shenzhen, 518048, China}

\author{Qi-Ming Ding}
\affiliation{Center on Frontiers of Computing Studies, Peking University, Beijing 100871, China}
\affiliation{School of Computer Science, Peking University, Beijing 100871, China}

\author{Haozhao Wu}
\affiliation{International Quantum Academy, Shenzhen, 518048, China}
\affiliation{College of Computer Science \& Software Engineering, Shenzhen University, Shenzhen, China}

\author{Mei-Shi Su}
\affiliation{Southern University of Science and Technology, Shenzhen, 518055, China}
\affiliation{International Quantum Academy, Shenzhen, 518048, China}

\author{Ya-Li Mao}
\affiliation{School of Physics, Nankai University, Tianjin, 300071, China}

\author{Xiao Yuan}
\email{xiaoyuan@pku.edu.cn}
\affiliation{Center on Frontiers of Computing Studies, Peking University, Beijing 100871, China}
\affiliation{School of Computer Science, Peking University, Beijing 100871, China}

\author{Zheng-Da Li}
\email{lizhengda@iqasz.cn}
\affiliation{International Quantum Academy, Shenzhen, 518048, China}


\begin{abstract}
\pagewiselinenumbers
\modulolinenumbers[1]
\linenumbers
Estimating the properties of quantum states is a fundamental task in both theoretical and experimental quantum physics. While numerous Pauli-based measurement strategies have been proposed, none explicitly exploit prior knowledge of structural features such as symmetry. Here, we introduce a compact measurement scheme that exploits symmetry to speed up quantum measurements. By incorporating symmetry into the measurement design, \red{our approach can rigorously reduce sample complexity for symmetric states and suitable observables.} We provide a theoretical analysis showing that, for specific observables, the variance of our estimator admits polynomial or even exponential improvements in the required number of state copies. To demonstrate practicality, we implement the scheme on photonic quantum platforms, preparing high-fidelity Greenberger-Horne-Zeilinger (GHZ) and $W$ states. The experimental results confirm substantial gains in estimation accuracy for both linear and nonlinear properties. Our work suggests that symmetry and other structural properties of quantum states can be systematically exploited as algorithmic resources, opening new avenues for scalable characterization and verification of complex quantum systems on near-term and fault-tolerant devices.

%our approach rigorously reduces sample complexity and generalizes existing Pauli-based methods.

 %   We experimentally prepare two kinds of high-quality three- and four-qubit symmetric states using an optical setup, respectively. 
\end{abstract}

\maketitle
\pagewiselinenumbers
\modulolinenumbers[1]
\linenumbers



%\section{Introduction}
\emph{Introduction ---}
Measuring linear and nonlinear properties of quantum states is a fundamental task in quantum computing and quantum information science. Linear properties of quantum states $\rho$ correspond to expectation values $\mathrm{Tr}(\rho H)$ of observable $H$, which constitute the key outputs of many quantum algorithms, including quantum simulation~\cite{lloyd1996universal,endo2020variational}, quantum machine learning~\cite{Schuld2015IntroQML, Biamonte2017QuantumML, schuld2019quantum}, quantum certification schemes~\cite{eisert2020quantum}, etc.
Beyond linear properties, {nonlinear} functionals of quantum states~\cite{ekert2002direct}, $\mathrm{Tr}(\rho^k H)$ for $k=2,3,\dots$, also play a central role in various contexts, including the measurement of R\'enyi entropies~\cite{islam2015measuring} and out-of-time-ordered correlators~\cite{swingle2016measuring}. 
Since the observable $H$ may contain many Pauli terms, estimating these quantities typically requires many state copies, reducing the number of measurements is therefore critical.
%Since the observable $H$ could be very complicated, the estimation of these quantities typically requires multiple copies of the quantum state, reducing the number of required measurements is of critical importance.

% Many quantum algorithms rely on estimating \emph{linear} properties of a quantum state~$\rho$, such as the expectation value $\mathrm{Tr}(\rho H)$ of an observable~$H$. Representative examples include the Harrow–Hassidim–Lloyd (HHL) algorithm~\cite{harrow2009quantum}, quantum machine learning protocols~\cite{schuld2019quantum}, and quantum certification schemes~\cite{eisert2020quantum}. 




%existing measurement algorithms



% The central idea is to estimate the expectation value of an observable $H = \sum_j \alpha_j Q_j$ where $\alpha_j$ is real number and $Q_j$ is a Pauli operator, for a state $\rho$ by employing optimized grouping strategies—guided by various loss functions—to measure as many Pauli bases $Q_j$ as possible within each setting.

%leverage state properties

Following the seminal work on classical shadow (CS) tomography~\cite{huang2020predicting}, 
which enables the estimation of expectation values of many observables, 
numerous randomized measurement schemes based on Pauli-basis measurements have been proposed~\cite{Huang2021Efficient,hadfield2022measurements,wu2023overlapped,Shlosberg2023Adaptive,zhang2023composite,gresch2025guaranteed}. 
These methods avoid the use of two-qubit gates before measurement and are therefore practical for near-term quantum devices~\cite{huang2020predicting}. 
In many tasks, partial information about the quantum state is often available and can, in principle, be exploited to reduce measurement costs. \red{For example, Tóth et al.~\cite{toth2010permuted} leveraged permutation invariance to efficiently reconstruct the permutationally invariant (PI) component of a quantum state.}
Such prior knowledge may include symmetry under a group action~\cite{toth2010permuted,elben2020mixed}, 
known conserved quantities such as particle number or total spin~\cite{bravyi2017tapering,setia2020reducing}, 
sparsity in the Pauli decomposition~\cite{peruzzo2014variational}, 
low-rank or pure-state structure~\cite{gross2010quantum}, 
or restrictions to a specific subspace, such as the symmetric subspace~\cite{stockton2003characterizing,moroder2012permutationally}. 
Nevertheless, existing randomized measurement approaches~\cite{huang2020predicting,Huang2021Efficient,hadfield2022measurements,wu2023overlapped,Shlosberg2023Adaptive,zhang2023composite,gresch2025guaranteed} 
rarely take advantage of such structural information.

In particular, symmetry in quantum states is of particular importance, with broad implications in quantum error correction~\cite{ouyang2014permutation,ouyang2017permutation}, quantum communication~\cite{dur2000three}, and photonic quantum computing~\cite{tichy2015sampling}. 
By exploiting permutation symmetry, 
Zhang and Tong proposed an optimal entanglement measurement method for measuring entanglement that exploits symmetric properties of quantum states in measurement~\cite{zhang2024inferring}. While this approach achieves optimal performance in theory, it requires implementing deep quantum circuits, which poses challenges for near-term quantum devices.




%propose our algs
Here, we propose a compact randomized Pauli-basis measurement scheme tailored for symmetric quantum states by transferring the symmetry from states to observables. 
\red{This leads to a measurement-compatible symmetry-averaging protocol that preserves the underlying Pauli measurement structure while reducing redundant measurement contributions. We prove that, for any fixed Pauli-based measurement protocol, the resulting estimator has no larger variance than the original estimator. In particular, for full-rank symmetric states, this holds whenever the averaging changes at least one measurable block.} 
%enhances the efficiency of quantum measurements and outperforms existing Pauli-based schemes~\cite{huang2020predicting,Huang2021Efficient,hadfield2022measurements,wu2023overlapped,Shlosberg2023Adaptive,zhang2023composite,gresch2025guaranteed}, while avoiding the overhead associated with the deep quantum circuits required in Ref.~\cite{zhang2024inferring}.
Experimentally, we estimate the expectation values of a spin Hamiltonian and an randomly generated Hermitian observable with respect to GHZ and $W$ states on photonic quantum devices. We further apply the method to evaluate nonlinear functionals, $\tr{\rho^2 H}$. Experimental results consistently show superior performance over existing Pauli-basis measurement schemes~\cite{huang2020predicting,Huang2021Efficient,hadfield2022measurements,wu2023overlapped,Shlosberg2023Adaptive,zhang2023composite,gresch2025guaranteed}.

\begin{figure}[t]
    \centering
    \includegraphics[width=0.85\linewidth]{figs/fig-theo.pdf}
 \caption{Illustration of the compact measurement scheme. For a permutation-symmetric state $\rho$, estimating $\alpha_i Q_i$ is equivalent to estimating its twirled counterpart $\mathcal{S}_{\mathcal{X}}(\alpha_i Q_i)$.}
    \label{fig:theo}
\end{figure}

%%%introduce theory
\emph{Compact measurement scheme ---}
For simplicity, we focus on estimating linear properties in the main text, \red{the nonlinear extension is summarized later and detailed in Supplementary Information (SI) Sec.~III.}
%and refer to Supplementary Information (SI) for the results of nonlinear properties.


Consider a symmetric quantum state $\rho$, which is given in a quantum oracle. Suppose the symmetry is described by a finite or compact Lie group $G$ satisfying $U_g \rho U_g^\dagger = \rho$, where $g\in G$ and $U_g$ is the unitary representation of $g$. 
Then $\rho$ is invariant under the twirling channel 
$    \Scal(\cdot) = \frac{1}{\abs{G}} \sum_{g\in G} U_g (\cdot) U_g^{\dagger},
$
for finite group $G$ with the number of elements $\abs{G}$, and $\Scal(\cdot) = \int_{g} d\mu(g) U_g (\cdot) U_g^\dagger$ for a compact Lie group $G$ with the Haar measure $\mu(g)$.


For observable $H$, suppose it is expanded in the Pauli basis as $H=\sum_{j=1}^m \alpha_j Q_j$ with $m$ non-trivial terms.
Since $\mathcal{S}(\rho) = \rho$, the expectation value satisfies 
$\mathrm{Tr}(\rho H) = \mathrm{Tr}(\mathcal{S}(\rho) H)$, 
and thus $\rho$ can be replaced by $\mathcal{S}(\rho)$ without changing the estimated value. 
As implementing the transformation $U_g$ on $\rho$ is costly, 
we instead apply $U_g$ to the observable $H$, using the cyclic property of the trace $
    \mathrm{Tr}(\mathcal{S}(\rho) H) = \mathrm{Tr}(\rho\, \mathcal{S}(H))$.

% Suppose $T$ number of samples are allowed to estimate $\tr{\rho H}$. For a Pauli-based algorithm, by utilizing an optimization function $f(H = \sum_{j=1}^m Q_j, T)$, it generates 
% measurement $P^{(k)}$, with the frequency $n_k$ for $k\in [L]$, such that $\sum_{k = 1}^L n_k = T$.
% The estimator can be defined as
% \begin{equation}
%     v = \sum_{k = 1}^L \sum_{j=1}^m \alpha_{j} \sum_{t = 1}^{n_k}\frac{\mu(P^{(k)}|Q_j,t)}{n_k} 
% \end{equation}
% where $\mu(P^{(k)}|Q_j)$ is nonzero only when $Q_j \triangleright P^{(k)}$ holds, $Q_j \triangleright P^{(k)}$ means the $i$-th qubit of $Q_j$ either equals identity or equals the $i$-th qubit of $P^{(k)}$. When $Q_j \triangleright P^{(k)}$, $\mu(P^{(k)}|Q_j,t) = \prod_{i:Q_{j,i}\ne \Ibb} \mu(P^{(k)}_i,t)$, and $\mu(P^{(k)},t)_i\in \cbra{\pm 1}$ is the outcome of the $t$-th round and $i$-th qubit with measurement $P^{(k)}$.
Suppose that $T$ measurement shots are available to estimate $\mathrm{Tr}(\rho H)$, where $H=\sum_{j=1}^{m}\alpha_j Q_j$ is a Pauli-string decomposition of the observable. 
In a Pauli-based algorithm, an optimization rule $f(H,T)$ selects Pauli measurement settings $\{P^{(k)}\}_{k=1}^{L}$ with corresponding shot counts $\{n_k\}_{k=1}^{L}$ satisfying $\sum_{k=1}^{L}n_k=T$. 
\red{We use the convention that $\mu(P^{(k)}\mid Q_j,t)=0$ whenever $Q_j\not\triangleright P^{(k)}$, i.e., whenever the Pauli string $Q_j$ cannot be inferred from measurements in the basis $P^{(k)}$. 
With this convention, the estimator is}
\begin{equation}
    v(H,f)
    =
    \sum_{j=1}^{m}\frac{\alpha_j}{\red{\mathrm{Freq}(Q_j)}}
    \sum_{k=1}^{L}
    \sum_{t=1}^{n_k}
   \mu(P^{(k)}\mid Q_j,t) ,
    \label{eq:estimation_formula}
\end{equation}
\red{where $P^{(k)}$ denotes the $k$th Pauli measurement basis selected by $f$, $n_k$ is the number of shots assigned to this basis, $\mathrm{Freq}(Q_j):=\sum_{k:\,Q_j\triangleright P^{(k)}}n_k$ is the total number of shots that effectively measure $Q_j$, and $Q_j\triangleright P^{(k)}$ means that $Q_{j,i}$ is either $\mathbb I$ or equal to $P_i^{(k)}$ for every qubit index $i$. Here, we assume \(\operatorname{Freq}(Q_j)>0\) for all retained Pauli terms. When \(\operatorname{Freq}(Q_j)=0\), the corresponding term cannot be estimated from the selected measurements and may be omitted, leading to a biased estimator. Such an approximation is reasonable only when the associated coefficient \(|\alpha_j|\) is sufficiently small.
When $Q_j\triangleright P^{(k)}$, the single-shot estimator is $\mu(P^{(k)}\mid Q_j,t)=\prod_{i:\,Q_{j,i}\neq\mathbb I}\mu(P_i^{(k)},t)$, where $\mu(P_i^{(k)},t)\in\{\pm1\}$ is the measurement outcome on qubit $i$ in the $t$th shot. 
Thus, Eq.~\eqref{eq:estimation_formula} averages all valid single-shot estimates of each Pauli expectation value $\mathrm{Tr}(\rho Q_j)$ and then combines them using the coefficients $\alpha_j$.}

Two Pauli strings are qubit-wise commuting, or compatible, if they commute on each qubit individually. 
A compatible set is a set of Pauli strings that are mutually qubit-wise commuting. 
For example, for any fixed measurement basis $P^{(k)}$, all Pauli strings $Q_j$ satisfying $Q_j\triangleright P^{(k)}$ form a compatible set.




\begin{figure*}[t]
	\centering
	\includegraphics[width=0.9\linewidth]{figs/fig-exp.pdf}
	\caption{Experimental setup. The schematics are divided into three modules: (a) the state preparation module, (b) the measurement module, and (c) the data-processing module. We switch four kinds of symmetric states $\{|\text{GHZ}_3\rangle, |W_3\rangle, |\text{GHZ}_4\rangle, |W_4\rangle\}$ by adjusting specific angle combinations $\mathcal{A}=[\theta,\phi_1,\phi_2,\phi_3,\phi_4,\phi_5]$ of waveplates labeled in the figure. The measurement apparatus performs polarization and path measurements on the signal (idler) photon using a set of QWP-HWP pairs placed before and after the BD, respectively.
    }
	\label{fig-exp}
\end{figure*}

% For simplicity, we assume $G$ is a discrete group. 
% Given a Pauli operator $P$, consider a subset $\{U_{g_1}, \ldots, U_{g_N}\}$ with $g_i \in G$ such that
% \begin{equation}
%      \widetilde{P} = \frac{1}{N} \sum_{i=1}^N U_{g_i} P U_{g_i}^\dagger 
% \end{equation}
% has fewer than $N$ incompatible sets in its Pauli decomposition. 
% In this case, the estimator $v(\widetilde{P},f)$ achieves a strictly smaller variance than $v(P,f)$ for the same number of measurements.  


% For permutation-invariant states, we further show that the variance can be reduced below $\mathcal{O}(1)$ for certain observables.  

\red{Here, we assume that the chosen
twirling subset $\mathcal{X}\subseteq G$ is taken such that each $S_{\mathcal{X}}(Q_j)$ can be expressed as a finite sum of signed
Pauli strings. A more detailed analysis is provided in SI Sec.~I.A.} Since $\rho$ is invariant under all transformations $U_g$, we define a symmetric twirling operation as the average over a subset $\Xcal \subseteq G$, which symmetrizes each $Q_j$ into \red{$J$ distinct Pauli terms $\{W_i\}$, yielding}
\begin{equation}
    \Scal_{\Xcal}(Q_j) = \frac{1}{J}\sum_{i=1}^J \operatorname{sgn}(W_i)\,W_i,
\end{equation}
\red{where $\operatorname{sgn}(W_i) \in \{\pm 1\}$ denotes the sign of $W_i$. We refer to the set of Pauli strings appearing in 
\(\mathcal S_{\Xcal}(Q_j)\) as the symmetry orbit of \(Q_j\).} We illustrate this process in Fig.~\ref{fig:theo}. The compact grouping algorithm first expresses the observable $H$ in its symmetric-twirled form, and then uses a tailored grouping procedure together with Eq.~\eqref{eq:estimation_formula} to estimate the expectation value of $\mathcal{S}_{\mathcal{X}}(H)$.
\red{When the observable contains only a single Pauli operator, $H=Q$, the variance-reduction effect of symmetric twirling can be established directly. 
Specifically, for any fixed admissible estimator-construction rule $f$, applying the same rule to the twirled observable $\mathcal S_{\Xcal}(Q)$ gives $\mathrm{Var}\!\left(v(\mathcal S_{\Xcal}(Q),f)\right)\leq \mathrm{Var}\!\left(v(Q,f)\right)$. 
A detailed derivation is provided in SI Sec.~II.A.
Moreover, the inequality is strict whenever the decomposition of $\mathcal S_{\Xcal}(Q)$ contains at least two distinct jointly measurable Pauli strings $W_i$ and $W_j$ whose correlation on the state $\rho$ is not saturated, e.g., $\abs{\operatorname{Tr}(\rho W_iW_j)}<1$. 
Thus, for a single Pauli term, symmetric twirling cannot increase the estimator variance and can strictly reduce it when compatible Pauli components in the twirled orbit provide nontrivial averaging. 
Below, we extend this observation to a general Pauli observable $H=\sum_i \alpha_i Q_i$ by deriving a corresponding variance comparison bound.} 
 %This construction enables us to establish the following variance bound.
\begin{theorem}[Informal]
Let $\Xcal \subseteq G$ be a symmetric twirling set for all Pauli operators $\{Q_i\}$ in the decomposition $H=\sum_{i=1}^m \alpha_i Q_i$.  
Then there exists a simple compact grouping estimator $\widehat{v}$ for $\text{tr}(\rho H)$ that measures one basis per compatible set and whose variance obeys the informal bound
\begin{equation}
\mathrm{Var}(\widehat{v})
\;\lesssim\;
\sum_{i=1}^{m}\alpha_i^2\,\frac{\Im_i}{Y_i}\;+\;\mathcal{R}(H,\rho),
\end{equation}
 where $Y_i$ denotes the number of Pauli terms in the decomposition of $\Scal_{\Xcal}(Q_i)$, $\Im_i$ denotes the number of compatible sets among these terms, and $\mathcal{R}(H,\rho)$ is an intra-group covariance remainder (coming only from correlations of Pauli terms within the same measured basis).
\label{thm:main_symmetric_average}
\end{theorem}
The covariance remainder $\Rcal(H,\rho)$ quantifies the correlation between measurement outcomes of different Pauli operators within the same group, and it is often close to zero for many choices of $\rho$ and $H$, especially when noise suppresses correlations in the quantum state while approximately preserving the relevant symmetry.
%when $\rho$ is strongly affected by noise.
We provide the formal theorem and its discussion in the \red{SI Sec.~II.A}. 
Here, we give some examples of the variance assuming $\Rcal(H,\rho)$ is close to zero. Consider the extreme case of a Pauli operator $P$ acting on $n$ qubits, specified by the triple $(n_x, n_y, n_z)$, where $n_x$, $n_y$, and $n_z$ denote the numbers of local Pauli-$X$, $Y$, and $Z$ operators, respectively, with $n \ge n_x + n_y + n_z$.
 After applying the twirl of the permutation group, the variance can be reduced from $\Ord{1}$ to ${1}/{{n - n_1 - n_2 \choose n_3}},$
where $n_1 \leq n_2 \leq n_3$ are the sorted values of $(n_x, n_y, n_z)$ if we assume $\mathcal{R}(H,\rho)\approx 0$. 
%In particular, \red{when $(n-n_1 - n_2)/n_3=cn$ for a constant $c\in(0,1)$,} 
%\red{i.e. $0<n_3/(n-n_1-n_2)<1$,} the scheme achieves an exponential reduction in sample complexity. 
\red{In particular, when $n-n_1-n_2=\Theta(n)$ and $n_3/(n-n_1-n_2)=c$ for some constant $c\in(0,1)$, the scheme achieves an exponential reduction in sample complexity. A detailed proof is provided in SI Sec.~II.B.
}
%$\binom{n-n_1-n_2}{n_3}$ grows exponentially with $n$. Therefore, 


%By contrast, without exploiting symmetry, \red{the variance remains $\Ocal(1)$.} 


The sample complexity can be further reduced when the Pauli terms in $\Scal_{\Xcal}(Q_i)$ across different $i$ are mutually compatible. As an example, consider the spin Hamiltonian
\begin{equation}
    H = \sum_{i\neq j} J_{ij} Z_i X_j ,
    \label{eq:spin_Ham}
\end{equation}
where $J_{ij}$ denote the coupling coefficients. 
Here we consider the case where $\tr{\rho X_j X_{j'}} = 0$ for all $j \neq j' \in [n]$, with $n$ the total number of qubits. In this situation, all corresponding correlation terms vanish, so the covariance remainder is zero for both the grouping and compact-grouping schemes. We note that the GHZ state with $n\geq 3$ fulfills this condition.
Define $e_i = \{Z_i X_j \,|\, j \neq i\}$ for $i \in [n]$, and measure each $e_i$ using the Pauli operator $P_i = Z_i X_{[n]\setminus \cbra{i}}$ with probability proportional to its weight $\abs{e_i} = \sum_{j:j\ne i} \abs{J_{ij}}$. In this setting, the estimation error scales as $
    \sum_{i,j:i\ne j} \abs{J_{ij}} \; \sum_{i} {\sum_{j:j\ne i} J_{ij}^2}/{\abs{e_i}}$.
By contrast, employing the compact grouping algorithm reduces the error to $({\sum_{i,j:i\ne j} J_{ij})^2}/{\,(n-1)\,}$,
under the assumption that the intra-group covariance remainder is negligible.
The derivations are provided in SI Sec.~II.A. Notably, when the couplings $J_{ij}$ are uniformly chosen from $\{-1,1\}$, the error scales as \red{$\Ord{n^3}$ in the original grouping scheme} and approaches zero for our compact scheme exploiting symmetry.

\red{
We note that the compact protocol may perform worse than an existing Pauli-based measurement protocol when \(\Rcal(H,\rho)\) is not negligible. Nevertheless, for any fixed Pauli-based protocol \(f\), we can construct a measurement-compatible symmetrized Hamiltonian by averaging complete symmetry orbits within each compatible set. This modification preserves the measurement protocol and is guaranteed to have no larger variance, as stated below.
}

\red{\begin{theorem}
Let \(H\) be a Hamiltonian, let \(\rho\) be a quantum state invariant under a symmetry group \(G\), and let \(f\) be a Pauli-based qubit-wise commuting measurement protocol. Then there exists a measurement-compatible symmetrized Hamiltonian \(H_{\mathrm{sym}}^{(f)}\), obtained by averaging complete symmetry orbits within each compatible set of \(f\), such that $ \mathrm{Tr} \left(\rho H_{\mathrm{sym}}^{(f)}\right)
    =
    \mathrm{Tr}(\rho H),$
and $ \mathrm{Var}\!\left[v(H_{\mathrm{sym}}^{(f)},f)\right]
    \le
    \mathrm{Var}\!\left[v(H,f)\right].$
Moreover, if \(\rho\) is full rank, the inequality is strict whenever \(H_{\mathrm{sym}}^{(f)}\neq H\).
\label{thm:symmetric_average_advantages}
\end{theorem}}
\red{The theorem mainly establishes a theoretical no-worse guarantee relative to any fixed Pauli-based measurement protocol. However, in certain cases, this measurement-compatible symmetrization may be less efficient than the compact protocol, which can exploit more of the symmetry information contained in the group $G$.
We provide a formal theorem and a more detailed proof in SI Sec.~II.C.}

\emph{Experimental setup ---} 
To demonstrate the advantage and practical feasibility of our method, we experimentally estimate both linear and nonlinear properties of some specific symmetric states, including three- and four-qubit GHZ \cite{greenberger1989} and $W$ states \cite{dur2000three,aguilar2015}, on a photonic platform. The experimental setup is illustrated in Fig.~\ref{fig-exp}, consisting of three main modules: state preparation, measurement, and data processing. 

\red{In the state-preparation module, photon pairs are first generated through the spontaneous parametric down-conversion (SPDC) process as photonic carriers, with quantum states encoded in their polarization and path degrees of freedom (DOFs). The  target state is controlled by adjustable angles  $\mathcal{A}=[\theta,\phi_1,\phi_2,\phi_3,\phi_4,\phi_5]$ of six half-wave plates (HWPs). Specifically, the settings $\mathcal{A}=[\frac{\pi}{8},-,\frac{\pi}{4},\frac{\pi}{4},0,0]$, $[\frac{1}{2}\arctan\sqrt{2},\frac{\pi}{4},\frac{\pi}{8},0,0,0]$, $[\frac{\pi}{8},-,\frac{\pi}{4},\frac{\pi}{4},\frac{\pi}{4},\frac{\pi}{4}]$, and $[\frac{\pi}{8},\frac{\pi}{4},\frac{\pi}{8},0,\frac{\pi}{8},0]$ correspond to the preparation of $|\mathrm{GHZ}_3\rangle$, $|W_3\rangle$, $|\mathrm{GHZ}_4\rangle$, and $|W_4\rangle$, respectively. The prepared states are then measured by Pauli measurements in the measurement module, with photon coincidence events recorded and analyzed in the data-processing module. Details are in the End Matter.}

\begin{figure*}[htbp!]
	\centering
\includegraphics[width =\textwidth]{figs/fig-data.pdf}
\caption{Estimation errors for properties of $W$ and GHZ states.  Expectation-value errors for a randomly generated three-qubit Hamiltonian with respect to the (a) $W$ state and (b) GHZ state; expectation-value errors for a four-qubit spin Hamiltonian with the (c) $W$ state and (d) GHZ state. Estimation errors of $\text{tr}(\rho^2 H)$ with three-qubit $\rho$ taken as the (e) $W$ state and (f) GHZ state, with $H$ a spin Hamiltonian. Results are shown for our algorithm, the Shadow Grouping method \cite{gresch2025guaranteed}, the derandomized approach \cite{Huang2021Efficient}, the overlapped grouping scheme \cite{wu2023overlapped}, and the adaptive Pauli measurement (AP) \cite{Shlosberg2023Adaptive}; AP is omitted in panels (e–f) due to its inferior accuracy and substantially higher runtime on this task. The estimation error in each point is calculated from the results of 20 independent experiments.}
%Estimation errors of $\text{tr}(\rho^2 H)$ for $\rho$ chosen from three-qubit (a) GHZ state and (b) $W$ state, with $H$ a spin Hamiltonian. Results are shown for our algorithm, the Shadow Grouping method \cite{gresch2025guaranteed}, the derandomized approach \cite{Huang2021Efficient}, and the overlapped grouping scheme \cite{wu2023overlapped}. SG exhibits inconsistent performance for certain measurement counts, likely due to specific features of the symmetric states examined.}
\label{fig:properties_numerics}
\end{figure*}



\textit{Linear properties ---}
Estimating linear properties of quantum states is central to quantum algorithm design and foundational studies~\cite{harrow2009quantum,gross2010quantum,kandala2017hardware,huang2020predicting}. We experimentally evaluate permutation-symmetric states ($W$ and GHZ) using our compact measurement scheme to estimate expectation values of benchmark Hamiltonians.

We consider two Hamiltonian families. First, the four-qubit spin model in Eq.~\eqref{eq:spin_Ham} with two-body couplings sampled uniformly from $[-1,1]$. Second, a randomly generated three-qubit Hamiltonian with coefficients drawn from $[-1/2^n,\,1/2^n]$ (with $n=3$), used to assess robustness beyond structured instances. The permutation group $\Tcal$ is adopted as the symmetric twirling set; we construct $\Scal_{\Tcal}(H)$ and determine measurement settings via overlapped grouping~\cite{wu2023overlapped}. The grouping subroutine is modular, \red{i.e., it acts as an interchangeable component that can be replaced by any other Pauli-based scheme without altering the overall framework.}

We benchmark against shadow grouping (SG)~\cite{gresch2025guaranteed}, derandomized classical shadows (Derand)~\cite{Huang2021Efficient}, overlapped grouping (OGM)~\cite{wu2023overlapped}, and adaptive Pauli scheme (AP)~\cite{Shlosberg2023Adaptive}. Each experiment is repeated $R=20$ times; the reported statistical error is the root-mean-square deviation $\epsilon \;=\; \sqrt{\frac{1}{R}\sum_{r=1}^{R}\!\bigl|\,v_r - \mathrm{Tr}(\rho H)\,\bigr|^{2}},
$ where $v_r$ denotes the estimator from the $r$-th repetition. 

\red{
Here the reference value is the ideal symmetric expectation value, rather than 
the expectation value of the reconstructed noisy experimental state. This choice 
matches the objective of the benchmark: estimating the target property of the 
ideal $W$ or GHZ state. For an imperfect experimental state 
$\hat{\rho}_{\rm exp}$, the compact estimator converges to 
$\operatorname{Tr}(\hat{\rho}_{\rm exp}H_{\rm sym})$ with 
$H_{\rm sym}=\mathcal{S}_{\mathcal T}(H)$, which is generally different from 
$\operatorname{Tr}(\hat{\rho}_{\rm exp}H)$. Equivalently, this is the expectation 
of $H$ on the symmetry-projected state 
$\mathcal{S}_{\mathcal T}(\hat{\rho}_{\rm exp})$. Thus, enforcing symmetry may 
produce a deterministic bias relative to the unsymmetrized noisy experimental 
value, but it can also suppress symmetry-breaking preparation errors and bring 
the estimate closer to the ideal symmetric target. We discuss this distinction 
in detail and provide an additional benchmark with respect to the noisy 
experimental reference in SI Sec.~IV.I.
}

Figures~\ref{fig:properties_numerics}(a–b) report the estimation of $\mathrm{Tr}(\rho H)$ for the random three-qubit Hamiltonian on the $W$ state (a) and the GHZ state (b), where the compact estimator exhibits a consistent advantage. Figures~\ref{fig:properties_numerics}(c–d) repeat the comparison for the spin model of Eq.~\eqref{eq:spin_Ham} on $W$ (c) and GHZ (d), highlighting pronounced gains on structured, physically motivated instances. These results indicate practical benefits of the proposed scheme for symmetric states on current quantum hardware.


%%%nonlinear properties
\textit{Nonlinear properties --- }
We also investigate the nonlinear properties of symmetric states, which play a central role in characterizing quantum correlations and entanglement beyond linear observables. In particular, nonlinear functionals such as $\tr{\rho^2 H}$ naturally arise in entanglement detection~\cite{guhne2009entanglement,cao2023generation}, purity estimation~\cite{ekert2002direct,brydges2019probing,wu2025state}, and the study of higher-order response functions~\cite{hamm2005principles}. As a representative example, we evaluate $\tr{\rho^2 H}$. In Fig.~\ref{fig:properties_numerics} (e)--(f), we benchmark our method against shadow grouping (SG), the derandomized classical shadows approach (Derand)~\cite{Huang2021Efficient}, and overlapped grouping measurement (OGM) on the $W$ state (e) and the GHZ state (f). We do not include the adaptive Pauli scheme~\cite{Shlosberg2023Adaptive}, which is substantially slower and does not outperform SG or OGM for the instances considered. For this comparison, we use a randomly generated three-qubit spin Hamiltonian of the form in Eq.~\eqref{eq:spin_Ham} with two-body couplings drawn uniformly from $[-1,1]$. Within our framework, $\tr{\rho^{2}H}$ admits a compact estimator by rewriting it as a quadratic form in single-copy Pauli expectations,
 $\tr{\rho^{2}H} = \sum_{i,j}\alpha_{ij}\tr{\rho P_i}\tr{\rho P_j}$ where $P_i,P_j$ are Pauli strings and $\alpha_{ij}$ are coefficients determined by $H$ and the chosen decomposition. Further implementation details are provided in SI Sec.~III. The experimental results show superior estimation precision for our method. These results demonstrate that our compact measurement scheme provides a robust and broadly applicable framework for estimating both linear and nonlinear observables of symmetric quantum states with significantly improved efficiency.


%%discussion
\textit{Conclusion and discussion --- }
In this work, we have introduced a compact measurement scheme tailored for quantum states with symmetry properties. We provide a theoretical analysis showing that, for specific observables, our method achieves substantial advantages, including an exponential reduction in the number of required state copies in extreme cases. We further demonstrated the practicality of the scheme through an experimental implementation on an optical quantum platform. Using high-fidelity GHZ and $W$ states, we observed clear improvements over existing Pauli-based measurement strategies, underscoring the utility of our approach for near-term quantum tasks. \red{We also provide a variance comparison between the Compact algorithm and existing measurement algorithms in SI Sec.~V, together with simulations of expectation-value estimation for the spin Hamiltonian and random Hamiltonian instances with respect to ideal GHZ and W states in SI Sec.~IV.H. 
These numerical results are consistent with the behavior shown in Fig.~\ref{fig:properties_numerics}.}

Our results highlight the broader significance of exploiting structural properties of quantum states to optimize measurement resources. An interesting open question is whether features beyond symmetry—such as entanglement structure or other conserved quantities—could be leveraged to achieve similar or even greater reductions in sample complexity while maintaining high estimation accuracy.\\

\emph{Acknowledgments ---} 
This work is supported by the Quantum Science and Technology-National Science and Technology Major Project (Nos.~2023ZD0301200,~2023ZD0300200),
the National Natural Science Foundation of China Grant (Nos.~92476103,~62375117,~T2522017,~12574394,~12361161602, and~12405014) and NSAF Grant (No.~U2330201), the Guangdong Basic and Applied Basic Research Foundation Grant (No.~2025B1515020064), Beijing Natural Science Foundation Grant (Nos.~1254053, and~Z250004), Beijing Science and Technology Planning Project Grant (No.~Z25110100810000), and the High-performance Computing Platform of Peking University.

\bibliography{ref}

\red{\appendix{\emph{End Matter ---}} We begin with the experimental state-preparation module, which provides a flexible photonic platform for generating the input states used throughout this work. As shown in Fig.~\ref{fig-exp}(a), this module can efficiently generate three- and four-qubit GHZ and $W$ states, encoded on the polarization and path DOFs of a pair of photons, by adjusting the angles of a set of HWPs. Here we employ a type-0 SPDC process that occurs in a periodically poled MgO-doped lithium niobate (PPLN) crystal embedded within a Sagnac interferometer to generate an entangled photon pair described by the state $|\psi\rangle=\cos 2\theta |H_sH_i\rangle+\sin 2\theta |V_sV_i\rangle$~\cite{li2022,mao2023}, where $|H\rangle$ and $|V\rangle$ represent the horizontal and vertical polarization states, respectively, and the subscripts $s$ and $i$ indicate the signal and idler photons, respectively. %More details about the two-photon entanglement source are shown in SI.


%We encode the desired states using the polarization and path degrees of freedom (d.o.f.) of the photon pairs. 
%The three-qubit system is composed of the polarization and path d.o.f.~of one photon and the polarization d.o.f.~of the other photon, while the four-qubit system additionally incorporates the path d.o.f.~of the second photon. Here, the horizontal polarization state $|H\rangle$ (upper path $|h\rangle$) and vertical polarization state $|V\rangle$ (lower path $|v\rangle$) represent $|0\rangle$ and $|1\rangle$, respectively.



%The state-preparation module in Fig.~\ref{fig-exp}(b) is designed to further prepare three- and four-qubit four-qubit GHZ and $W$ states. 

Further, we introduce the path DOF using a group of beam displacers (BDs) alternated with adjustable HWPs, where $|u\rangle$ and $|l\rangle$ represent the upper and lower path states of photons, respectively.  Here, the BDs are constructed by bonding polarizing beam splitters (PBSs) and mirrors. Including an HWP $H_0$ acting on the pump photons before the SPDC process, by adjusting the angles $\mathcal{A}=[\theta,\phi_1,\phi_2,\phi_3,\phi_4,\phi_5]$ of a specific combination of adjustable HWPs marked in Fig. \ref{fig-exp} (a), we can generate four multi-qubit symmetric states, including three- and four-qubit GHZ and $W$ states. By choosing $\mathcal{A}=[\frac{\pi}{8},-,\frac{\pi}{4},\frac{\pi}{4},0,0]$, we generate the state $|\text{GHZ}_3\rangle=\frac{1}{\sqrt{2}}(|H_su_sH_i\rangle+|V_sl_sV_i\rangle)$, where the minus sign indicates that the corresponding HWP is absent. Similarly, using $\mathcal{A}=[\red{\frac{1}{2}\arctan{\sqrt{2}}},\frac{\pi}{4},\frac{\pi}{8},0,0,0]$, we achieve the state $|W_3\rangle=\frac{1}{\sqrt{3}}(|H_su_sV_i\rangle+|H_sl_sH_i\rangle+|V_su_sH_i\rangle)$. Note that the prepared three-qubit quantum states do not rely on the path DOF of the idler photon; rather, they are achieved by counteracting the effects of the two BDs on the idler photons' paths through the configurations of the HWPs. By selecting $\mathcal{A}=[\frac{\pi}{8},-,\frac{\pi}{4},\frac{\pi}{4},\frac{\pi}{4},\frac{\pi}{4}]$, we obtain the four-qubit GHZ state $|\text{GHZ}_4\rangle=\frac{1}{\sqrt{2}}(|H_su_sH_iu_i\rangle+|V_sl_sV_il_i\rangle)$. Furthermore, the four-qubit $W$ state $|W_4\rangle=\frac{1}{2}(|H_su_sH_il_i\rangle+|H_su_sV_iu_i\rangle+|H_sl_sH_iu_i\rangle+|V_su_sH_iu_i\rangle)$ is prepared when $\mathcal{A}=[\frac{\pi}{8},\frac{\pi}{4},\frac{\pi}{8},0,\frac{\pi}{8},0]$.

Next, the prepared multi-qubit symmetric states are directed into a measurement module as depicted in Fig.~\ref{fig-exp}(b). For each incoming photon, this module sequentially performs Pauli measurements on the polarization qubit and the path qubit encoded within it. The measurement settings for the polarization qubit of the signal (idler) photon are controlled by the quarter-wave plate (QWP) and an HWP  before the BD, while the path qubit measurement settings are governed by the QWP and HWP behind it. Notably, in the three-qubit cases, the path measurement of the idler photon is not performed. Both the measurement outcomes of the polarization and path qubits are obtained from response events across the superconducting nanowire single-photon detectors (SNSPD). As shown in Fig.~\ref{fig-exp}(c), a data processing module, comprising a computer and a time-to-digital converter (TDC), is employed to chronologically record the photon coincidence events and analyze the experimental data for estimating the properties of the quantum states. Further experimental details are provided in SI Sec.~IV.}


\end{document}
